\subsection{Affine schemes}

In this section we briefly review the definition of an affine scheme.

\bdefn

    Let $A$ be a commutative ring; its {\emphcolor spectrum} is the topological space whose underlying set is the set of all prime ideals of $A$.
    We denote it by $\spec(A)$, and for a prime ideal $\p\leq A$ we denote the element of $\spec(A)$ by $[\p]\in\spec(A)$.

    For a subset $S\subseteq A$ we define its {\emphcolor vanishing set} to be
    $$ \vof(S) = \set{[\p]\in\spec(A)}[S\subseteq\p] . $$
    The collection of all such vanishing sets, $\set{\vof(S)}_{S\subseteq A}$, forms a topology of closed sets on $\spec(A)$.
    This topology is called the {\emphcolor Zariski topology}.

\edefn

Notice that for a subset $S\subseteq A$, $\vof(S)=\vof((S))$ where $(S)\leq A$ is the ideal generated by $A$.
Furthermore,
$$ \vof(\emptyset) = \vof(0) = \spec(A),\qquad \vof(A) = \emptyset,\qquad \bigcap_{i\in I}\vof(S_i) = \vof\parens{\bigcup_{i\in I}S_i} , $$
and that for two ideals $I,J\leq A$ we have that $\vof(I)\cup\vof(J)=\vof(IJ)$.
It follows that the Zariski topology is indeed a topology.
Notice that $\vof$ is order-reversing: if $S_1\subseteq S_2$ then $\vof(S_2)\subseteq\vof(S_1)$.

Recall that if $\phi\colon A\to B$ is a ring morphism then it pulls back prime ideals: if $\q\leq B$ is prime, then so is $\phi^{-1}\q\leq A$.
Thus $\spec$ is functorial:
$$ \spec\colon\Ring^\op\to\Top,\qquad \phi\mapsto\phi^* $$
where $\phi^*\colon\spec(B)\to\spec(A)$ is $[\q]\mapsto[\phi^{-1}\q]$.
This is continuous: for $I\leq A$
$$ (\phi^*)^{-1}\vof_A(I) = \set{[\q]\in\spec(B)}[I\subseteq\phi^{-1}\q] = \set{[\q]\in\spec(B)}[\phi(I)\subseteq\q] = \vof_B(\phi(I)) . $$

If $I\leq A$ is an ideal then its {\emphcolor radical} is the ideal
$$ \sqrt I = \set{a\in A}[\hbox{$a^n\in I$ for some $n\geq0$}] \leq A . $$
It is easy to see that $\vof(I)=\vof(\sqrt I)$: if $I\leq\p$ then if $f\in\sqrt I$ then $f^n\in I\leq\p$ for some $n$ and since $\p$ is prime, $f\in\p$.
Furthermore it can be seen that $\sqrt I$ is the intersection of all prime ideals containing $I$:
$$ \sqrt I = \bigcap_{I\leq\p}\p = \bigcap\vof(I) . $$
Clearly the left-hand side is contained in the right.
Conversely if $f\notin\sqrt I$ then $\sqrt I$ and $\set{f^n}_{n\geq0}$ do not intersect, and so $(\sqrt I)_f\leq A_f$ is a proper ideal.
Thus it is contained in a prime ideal $(\sqrt I)_f\leq\p_f\leq A_f$ associated to a prime ideal $\p\leq A$, so that $\sqrt I\leq\p$.
But $f\notin\p$, so $f\notin\bigcap\vof(I)$.

We also define the {\emphcolor nilradical} of $A$ be the radical of the zero ideal:
$$ \nil(A) = \sqrt0 = \set{a\in A}[\hbox{$a^n=0$ for some $n\geq0$}] . $$
Thus we have that
$$ \nil(A) = \bigcap_{[\p]\in\spec(A)}\p . $$

We can also form a base for the Zariski topology of $\spec(A)$.
For $S\subseteq A$ let
$$ \dof(S) = \vof(S)^c = \set{[\p]\in\spec(A)}[S\not\subseteq\p] $$
form the open sets of the Zariski topology.
For $f\in A$ let
$$ \dof(f) = \dof\set f = \set{[\p]\in\spec(A)}[f\notin\p] . $$
It is clear that
$$ \dof(S) = \bigcup_{f\in S}\dof(f) $$
so that the $\dof(f)$ form a base of open sets for the topology.

We can map ideals of $A$ to closed subsets of $\spec(A)$, we can also map subsets of $\spec(A)$ to ideals of $A$.
For $X\subseteq\spec(A)$ we define the ideal
$$ \idof(X) = \bigcap_{[\p]\in X}\p \leq A . $$
This is also clearly order-reversing: if $X_1\subseteq X_2$ then $\idof(X_2)\subseteq\idof(X_1)$.

Notice that for an ideal $I\leq A$,
$$ \idof(\vof(I)) = \bigcap_{I\leq\p}\p = \sqrt I . $$
And conversely for a subset $X\subseteq\spec(A)$,
$$ \vof(\idof(X)) = \cl X , $$
the closure of $X$ in the Zariski topology.
Indeed: if $X\subseteq\vof(I)$ then for $[\p]\in X$ we have $I\leq\p$ so that $I\leq\bigcap_{[\p]\in X}\p=\idof(X)$.
Thus $\vof(\idof(X))\leq\vof(I)$, so that $\vof(\idof(X))$ is the smallest closed subset containing $X$; i.e.\ $\cl X$.

Thus $\vof$ and $\idof$ form order-reversing bijections
$$ \set{\hbox{radical ideals of $A$}} \simarrow{\vof(\bul)}{\idof(\bul)} \set{\hbox{closed subsets of $\spec(A)$}} . $$

It follows that $\vof(I)\subseteq\vof(J)$ iff $\sqrt J\leq\sqrt I$: applying $\idof(\bul)$ gives us the left-to-right direction, and $\vof(\sqrt I)=\vof(I)$ gives the right-to-left direction.
In particular $\vof(I)=\vof(J)$ iff $\sqrt I=\sqrt J$.
This carries over to $\dof$ as well obviously: $\dof(I)\subseteq\dof(J)$ iff $\sqrt I\leq\sqrt J$.
In particular $\dof(f)\subseteq\dof(g)$ iff $\sqrt f\leq\sqrt g$ iff $f^n\in(g)$ for some $n\geq0$, or equivalently $g$ is invertible in $A_f$ ($g$ is invertible in $A_f$ iff $(g)_f$ is unital,
iff $(g)$ intersects the localizing set, iff $f^n\in(g)$ for some $n\geq0$).

Because $\vof(0)=\spec(A)$, $\vof(I)=\spec(A)$ iff $\sqrt I=\nil(A)$ iff $I\subseteq\nil(A)$.
In particular $\dof(f)=\emptyset$ iff $f\in\nil(A)$.
On the other hand, $\vof(I)=\emptyset$ for an ideal $I$ iff $I=A$, since every proper ideal is contained in a prime ideal.
Thus $\dof(f)=\spec(A)$ iff $f$ is a unit.

We now discuss some topological results regarding $\spec(A)$.

\bdefn

    A topological space $X$ is said to be
    \benum
        \item {\emphcolor quasicompact} if every open cover of $X$ has a finite subcover.
        Explicitly, if $X=\bigcup_{i\in I}U_i$ for $U_i\subseteq X$, then there are $i_1,\dots,i_n\in I$ such that $X=\bigcup_{j=1}^nU_{i_j}$.
        \item {\emphcolor connected} if $X$ cannot be written as the disjoint union of two nonempty open sets.
        Equivalently, $X$ has no clopen subsets (subsets which are simultaneously open and closed) other than $X$ and $\emptyset$.
        Otherwise $X$ is {\emphcolor disconnected}.
        \item {\emphcolor irreducible} if $X$ cannot be written as the union of two proper closed subsets.
        That is if $X=Y\cup Z$ for $Y,Z$ closed then either $X=Y$ or $X=Z$.
        Equivalently the intersection of any two nonempty open subsets is nonempty.
    \eenum

\edefn

\bnote

    The term ``quasicompact'' may be odd to you, as this is sometimes taken to be the definition of a ``compact'' space.
    For geometers, oftentimes a compact space is a space which is both quasicompact and Hausdorff.

\enote

\bremark

    A space $X$ is quasicompact iff every collection of closed sets with the finite intersection property has nonempty intersection.
    That is, if $\set{Z_i}_{i\in I}$ is a collection of closed subsets of $X$ such that $\bigcap_{i\in I_0}Z_i$ is nonempty for all finite $I_0\subseteq I$, then $\bigcap_{i\in I}Z_i\neq\emptyset$.

\eremark

\bremark

    It is immediate that a space $X$ is irreducible iff all nonempty open subsets of $X$ are dense.

\eremark

It is true that $\spec(A)$ is always quasicompact.
It is sufficient to consider basic closed sets.
If $\set{\vof(f_i)}_{i\in I}$ has the finite intersection property, then
$$ \bigcap_{i\in I}\vof(f_i) = \vof(\set{f_i}_{i\in I}) $$
and if this is empty, then $\ideal{f_i}[i\in I]=A$.
So there are $i_1,\dots,i_n\in I$ such that $1\in(f_{i_1},\dots,f_{i_n})$ so that
$$ \emptyset = \vof(f_{i_1},\dots,f_{i_n}) = \bigcap_{j=1}^n\vof(f_{i_j}) , $$
contradicting the finite intersection property.

\bdefn

    An {\emphcolor idempotent} in a ring $A$ is an element $f\in A$ such that $f^2=f$.
    A {\emphcolor nontrivial idempotent} is an idempotent which is not $0$ or $1$.
    A ring $A$ is {\emphcolor connected} if it has no nontrivial idempotents; otherwise it is {\emphcolor disconnected}.

\edefn

\bprop

    Let $A$ be a ring.
    Then the following conditions are equivalent:
    \benum
        \item $A$ is disconnected;
        \item $A$ is isomorphic to the product of two non-trivial rings;
        \item $\spec(A)$ is disconnected.
    \eenum
    In particular $\spec(A\times B)\cong\spec(A)\amalg\spec(B)$.

\eprop

\bproof

    Suppose (1), then if $f\in A$ is a nontrivial idempotent let $g=1-f$; this is also a nontrivial idempotent.
    The ideal $A^f=(f)$ has a ring structure whose addition and multiplication is inherited from $A$, its identity is $f$.
    We define $\phi\colon A\to A^f\times A^g$ by $a\mapsto(fa,ga)$.
    This is injective because $fa+ga=a$, and it is surjective because $fg=0$ so $\phi(fa+gb)=(fa,gb)$.
    Thus $A\cong A^f\times A^g$, and these are nontrivial because $f,g\neq0$; so (1) impies (2).

    Conversely if $A\cong B\times C$ then $(1,0),(0,1)\in A$ are nontrivial idempotents so $A$ is disconnected.
    Thus (2) implies (1).

    Now again suppose (2), so consider the ring $A\times B$.
    Then notice that $\spec(A\times B)=\dof(1,0)\sqcup\dof(0,1)$.
    Indeed:
    $$ \dof(1,0)\cup\dof(0,1) = \dof((1,0),(0,1)) = \dof(1) = \spec(A\times B),\qquad \dof(1,0)\cap\dof(0,1) = \dof((1,0)\cdot(0,1)) = \dof(0) = \emptyset . $$
    Thus (2) implies (3).

    Notice that the projection maps $A\times B\to A,B$ induce $\spec(A)\amalg\spec(B)\to\spec(A\times B)$.
    It is easy to see that this maps $\spec(A)$ to $\dof(1,0)$ and $\spec(B)$ to $\dof(1,0)$, showing that this is an isomorphism.

    Now assume (3), so $\spec(A)=\vof(I)\sqcup\vof(J)$ where $I,J$ are nonzero (so $\vof(I),\vof(J)$ are proper subsets).
    This means that $\vof(I+J)=\emptyset$ so that $I+J=A$ and $\vof(IJ)=\spec(A)$ so that $IJ=0$.
    Thus by the Chinese remainder theorem, $A\to A/I\times A/J$ is an isomorphism, proving (2).
    \qed

\eproof

\bprop

    $\spec(A)$ is irreducible iff $\nil(A)$ is prime.

\eprop

\bproof

    If $\nil(A)$ is not prime then there are $f,g\notin\nil(A)$ such that $fg\in\nil(A)$.
    So $\vof(f),\vof(g)\neq\spec(A)$ and $\spec(A)=\vof(fg)=\vof(f)\cup\vof(g)$; thus $\spec(A)$ is reducible as desired.

    Conversely if $\spec(A)$ is reducible then suppose $\spec(A)=\vof(I)\cup\vof(J)$ for proper subsets, so $I,J$ are not subsets of $\nil(A)$.
    But then $\vof(IJ)=\spec(A)$ so $IJ\leq\nil(A)$, so $\nil(A)$ is not prime.
    \qed

\eproof

\bcoro[name={domspecirred}]

    If $A$ is a domain then $\spec(A)$ is irreducible.

\ecoro

\bnote

    While it is true that $\spec\prod_{i=1}^nA_i\cong\coprod_{i=1}^n\spec A_i$ for finite products, this does not hold for infinite products.
    There is a canonical map $\coprod_{i\in I}\spec A_i\to\spec\prod_{i\in I}A_i$; and it is easy to see that this map is an embedding.
    But if all the $A_i$ are nonzero and $I$ is infinite, the left hand side is not compact (because $\set{\spec A_i}_{i\in I}$ does not have a finite subcover) and thus it cannot be
    an isomorphism.

\enote

\bdefn

    Let $X$ be a topological space.
    A {\emphcolor closed point} of $X$ is a point $x\in X$ such that $\cl{\set x}=\set x$.

\edefn

In a Hausdorff space, all points are closed.
For $\spec A$, we know that $\cl{[\p]}=\vof(\idof([\p]))=\vof(\p)$.
Thus $[\p]$ is a closed point iff $\p$ is not contained in any prime ideal other than itself; equivalently it is a maximal ideal.
That is: $[\p]\in\spec A$ is a closed point iff $\p\leq A$ is maximal.

\bdefn

    A subset of a topological space $Z\subseteq X$ is an {\emphcolor irreducible component} if it is irreducible and maximally so: it is not properly contained in an irreducible set.

\edefn

\bprop

    A subset $Z\subseteq X$ is irreducible iff $\cl Z\subseteq X$ is.

\eprop

\bproof

    If $Z$ is irreducible and $\cl Z=Y\cup W$ are closed, then $Z\subseteq Y\cup W$ so that $Z\subseteq Y$ or $Z\subseteq W$.
    Then $\cl Z=Y$ or $W$ because they are closed in $\cl Z$ and thus in $X$.
    Conversely if $\cl Z$ is irreducible then if $Z=Y\cup W$, then $\cl Z=\cl Y\cup\cl W$ so $\cl Z=\cl Y$ or $\cl W$.
    Suppose the former, so $Z=\cl Y\cap Z=Y$ since $Y\subseteq Z$ is closed.
    \qed

\eproof

\bprop

    Every point $p\in X$ is contained in a (not necessarily unique) irreducible component, and all irreducible components of $X$ are closed.

\eprop

\bproof

    The first claim is due to Zorn.
    Clearly $\set p$ is irreducible, so we can form the nonempty poset of irreducible sets containing $p$.
    For a chain of such irreducible sets $\set{Z_i}_{i\in I}$, we have that $Z=\bigcup_{i\in I}Z_i$ is irreducible.
    Indeed if $Z=Y\cup W$ are closed then suppose $Y\subsetneq Z$ then there is an $f\in Z_i$ such that $f\notin Z$.
    Because $Z_i$ is irreducible we must have $Z_i\subseteq W$.
    Then for all $j\geq i$ we have that $Z_j\subseteq W$ for the same reason; then $Z=\bigcup_{j\geq i}Z_j=W$ so that $Z$ is irreducible as desired.

    The second claim follows from the previous proposition: if $Z$ is an irreducible component then $\cl Z$ is irreducible as well and since $Z\subseteq\cl Z$
    by maximality $Z=\cl Z$.
    \qed

\eproof

\bdefn

    Let $x,y\in X$ be points in a topological space such that $x\in\cl{\set y}$.
    Then $x$ is called a {\emphcolor specialization} of $y$, and $y$ a {\emphcolor generalization} of $x$.

    A point $p\in X$ is a {\emphcolor generic point} if $\cl{\set p}=X$; i.e.\ if it is a generalization of all the points of $X$.

\edefn

In $\spec A$ because $\cl{[\p]}=\vof(\p)$, $[\p]$ is a generalization of $[\q]$ iff $[\q]\in\vof(\p)$ iff $\p\leq\q$.
And $[\p]$ is a generic point iff $\vof(\p)=\spec(A)$ iff $\p=\nil A$, so that $\spec A$ has a generic point iff $\nil A$ is prime iff $\spec A$ is irreducible.

\bdefn

    A poset $\Sigma$ has the {\emphcolor descending chain condition} ({\emphcolor DCC}) if any of the following two equivalent conditions are satisfied:
    \benum
        \item every subset $\Sigma'\subseteq\Sigma$ has a minimal element;
        \item every chain of elements $p_1\geq p_2\geq p_3\geq\cdots$ eventually stabilizes: $p_n=p_{n+1}=p_{n+2}=\cdots$ for some $n$ onward.
    \eenum
    Dually $\Sigma$ has the {\emphcolor ascending chain condition} ({\emphcolor ACC}) if $\Sigma^\op$ has the DCC.

\edefn

The equivalence of the conditions is clear: (1) implies (2) by defining $\Sigma'=\set{p_n}_n$ and taking the minimum.
Conversely if $\Sigma'\subseteq\Sigma$ has no minimal element, then we can form an infinite decreasing chain.

\bdefn

    A topological space $X$ is {\emphcolor Noetherian} if its poset of closed subsets has the DCC (equivalently its poset of open subsets has the ACC).

\edefn

There is a notion of {\emphcolor Noetherian induction}: if $\Sigma$ has DCC and we have some property $P$ (which one can view as a function $\Sigma\to\set{0,1}$)
such that $P(x)$ holds for all minimal elements of $\Sigma$, and for $y\in\Sigma$ non-minimal if $P(x)$ holds for all $x<y$ then $P(y)$.
Then $P(x)$ holds for all $x\in\Sigma$.
This is easy to prove: take $\Sigma'$ to be the set of all $x\in\Sigma$ for which $P(x)$ does not hold.

Using Noetherian induction one can prove the following.

\bprop

    If $X$ is a Noetherian topological space, then every nonempty closed subset $Z\subseteq X$ can be uniquely expressed as a finite union $Z=Z_1\cup\cdots\cup Z_n$
    where the $Z_i$ are all irreducible closed subsets, none containing any other.

\eprop

\bremark

    It is clear that if $A$ is a Noetherian ring (its poset of ideals has ACC) then $\spec A$ is a Noetherian space.
    The converse is not true: if $\spec A$ is a Noetherian space then the poset of $A$'s {\it radical\/} ideals has ACC.

\eremark

